\documentclass[11pt]{article}
\usepackage{amsmath,amsthm,amssymb}
\usepackage[margin=1in]{geometry}
\usepackage{booktabs}
\usepackage{enumitem}
\usepackage{ytableau}

\newtheorem{theorem}{Theorem}
\newtheorem{proposition}[theorem]{Proposition}
\newtheorem{lemma}[theorem]{Lemma}
\newtheorem{corollary}[theorem]{Corollary}
\newtheorem{conjecture}[theorem]{Conjecture}
\theoremstyle{definition}
\newtheorem{definition}[theorem]{Definition}
\newtheorem{remark}[theorem]{Remark}
\newtheorem{example}[theorem]{Example}

\newcommand{\Om}{\Omega}
\newcommand{\Omstar}{\Omega^{\!*}}
\newcommand{\la}{\lambda}
\newcommand{\hla}{\hat{\lambda}}
\newcommand{\SYT}{\mathrm{SYT}}
\newcommand{\NN}{\mathbb{N}}
\newcommand{\QQ}{\mathbb{Q}}
\newcommand{\ZZ}{\mathbb{Z}}

\title{The SR-companion conjecture is false}
\author{Clio}
\date{19 May 2026}

\begin{document}
\maketitle

\begin{abstract}
Yesterday we proved a \emph{universal SC necessity} theorem: for every $\la$
and every $T \in \SYT(\la)$ with $p_T(q) = 0$, the tableau $T$ has an SC
pair, where $p_T(q) = \Om^{(\la)}_{T,T}$ is the Hoefsmit diagonal of
$\Om^{(\la)} = R'_{\ell+1} \cdots R'_n$. The proof is a one-line application
of Hoefsmit positivity in the $b' = 1$ basis to the all-diagonal term of the
matrix product expansion.

The mirror conjecture was natural: applying the same rigidity to
$\Omstar = L_n \cdots L_{\ell+1}$ in the $b = 1$ rescaling should give
universal SR necessity. We report that this conjecture is \emph{false}.
Across a twelve-shape computational corpus with $514$ zero diagonals
total, $193$ have an SC pair but \emph{no} SR pair --- a $38\%$ refutation
rate. The cleanest counterexample is $\la = (4,3)$ at
$T = ((1,3,5,7), (2,4,6))$, where every adjacent pair $(i, i+1)$ is a
descent and no adjacent pair lies in the same row.

We localise where the symmetric argument fails: the rescaling from
$b' = 1$ to $b = 1$ does \emph{not} carry positivity to $\Omstar$ through
the chain product $L_n \cdots L_{\ell+1}$ in the form required for the
all-diagonal rigidity. Diagnostic~3 stands in its current form (SC or SR,
not both); the converse of Diagnostic~4 must seek a Hecke-algebraic
witness, not a symmetric positivity argument.
\end{abstract}

\section{Setup}

Fix $\la = (\hla, 1^q)$ with $\ell(\hla) \ge 2$ and every row
$\hla_r \ge 2$. Set $\ell = \ell(\hla)$, $n = |\la|$. Write
$\Om = \Om^{(\la)} = R'_{\ell+1} R'_{\ell+2} \cdots R'_n$, where
$R'_k = S_{k-1} S_{k-2} \cdots S_1$ and $S_i = T_i + 1$ in $H_q(S_n)$.
Acting on $V^\la$ in the Hoefsmit (seminormal) basis indexed by $\SYT(\la)$
with the $b' = 1$ convention, the operator $\Om$ has a matrix; the diagonal
entries are
\[
  p_T(q) \;:=\; \Om_{T,T} \;\in\; \QQ(q).
\]
Let $\Omstar = L_n L_{n-1} \cdots L_{\ell+1}$ where $L_k = S_1 S_2 \cdots S_{k-1}$,
the leftmost--rightmost mirror of $\Om$.

\begin{definition}\label{def:scsr}
A standard tableau $T$ has an \emph{SC pair} (``same-column'') at $i$ if
the entries $i$ and $i+1$ lie in the same column of $T$. It has an
\emph{SR pair} (``same-row'') at $i$ if $i$ and $i+1$ lie in the same row.
Equivalently, SC pairs are exactly the descents of $T$ (read as a word),
and SR pairs are the same-row ascents.
\end{definition}

\section{The conjecture and its motivation}

\begin{theorem}[Universal SC necessity, 2026-05-18]\label{thm:sc}
For every $\la$ with $\ell(\hla) \ge 2$, every row $\hla_r \ge 2$, and every
$T \in \SYT(\la)$, if $p_T(q) = 0$ then $T$ has an SC pair.
\end{theorem}

The proof has one ingredient: Hoefsmit positivity (Lemma~1 of
\texttt{2026-05-17-converse-per-syt-positivity.tex}) puts every entry of
$T_i + 1$ in $\QQ_{\ge 0}(q)$ in the $b' = 1$ basis. The matrix product
$\Om = \prod_j (T_{i_j} + 1)$ expands into a sum of nonneg-rational
terms indexed by paths through index sets. The diagonal entry $\Om_{T,T}$
is one such sum. The all-diagonal path contributes
$\prod_j (S_{i_j})_{T,T} = \prod_j [d_j(T) + 1]_q / [d_j(T)]_q$,
where $d_j(T)$ is the axial distance of $(i_j, i_j + 1)$ in $T$. Vanishing
of a sum of nonneg-rational terms forces every term to vanish; the
all-diagonal term vanishes iff some factor vanishes, i.e.\ some
$d_j(T) = -1$, i.e.\ $i_j$ is a column-descent of $T$. \hfill$\square$

The \emph{SR companion conjecture} was the apparent mirror:

\begin{conjecture}[SR companion --- now refuted]\label{conj:sr}
For every $\la$ and every $T \in \SYT(\la)$, if $p_T(q) = 0$ then $T$ has
an SR pair.
\end{conjecture}

The motivating argument: $\Omstar$ is the $\iota$-image of $\Om$, and
Diagonal~4 (\texttt{2026-05-18-d-class-row-zero.tex}) gives
$p_T = 0 \iff \Om v_T = 0$ or $\Omstar v_T = 0$. If $\Omstar$ were
nonneg-rational in some symmetric basis, the same all-diagonal rigidity
would give SR. The natural candidate basis is $b = 1$ (rescaling
$b' = 1$ Hoefsmit entries by axial-distance ratios). Then ``SR necessity''
would follow.

\section{Computational refutation}

We tested Conjecture~\ref{conj:sr} on the $12$-shape corpus used by
\texttt{dichotomy\_verifier.py}, evaluating at $q = 11$ (a generic
specialisation; matches the convention of
\texttt{fast\_numeric\_probe.py}). For each $T$ with $\Om_{T,T} = 0$, we
recorded whether $T$ has an SC pair, an SR pair, both, or neither.

\begin{table}[h]
\centering
\begin{tabular}{lrrrrrr}
\toprule
shape & $N$ & zeros & both & SC only & SR only & neither \\
\midrule
(3,2)         & 5   & 0   & 0   & 0   & 0 & 0 \\
(3,2,1)       & 16  & 4   & 0   & 4   & 0 & 0 \\
(2,2,1)       & 5   & 1   & 1   & 0   & 0 & 0 \\
(3,2,2)       & 21  & 11  & 9   & 2   & 0 & 0 \\
(3,2,1,1)     & 35  & 21  & 15  & 6   & 0 & 0 \\
(3,3,2)       & 42  & 26  & 17  & 9   & 0 & 0 \\
(3,2,2,1)     & 64  & 49  & 35  & 14  & 0 & 0 \\
(3,3,2,1)     & 168 & 134 & 92  & 42  & 0 & 0 \\
(4,3)         & 14  & 5   & 2   & 3   & 0 & 0 \\
(4,3,1)       & 70  & 36  & 21  & 15  & 0 & 0 \\
(4,3,2)       & 168 & 149 & 87  & 62  & 0 & 0 \\
(3,3,1,1,1)   & 90  & 78  & 42  & 36  & 0 & 0 \\
\midrule
\textbf{Total} & --- & \textbf{514} & \textbf{321} & \textbf{193} & 0 & 0 \\
\bottomrule
\end{tabular}
\caption{Distribution of $(\text{SC}, \text{SR})$ pair existence across
$514$ zero diagonals on a twelve-shape corpus. ``SR only = 0'' confirms
Theorem~\ref{thm:sc} (universal SC necessity). ``SC only = 193''
refutes Conjecture~\ref{conj:sr}. ``Neither = 0'' confirms Diagnostic~3
(SC or SR) holds throughout the corpus.}
\label{tab:results}
\end{table}

\begin{proposition}
Conjecture~\ref{conj:sr} is false. At least $193$ counterexamples exist
on shapes of size $\le 9$.
\end{proposition}

\section{The cleanest counterexample}

\begin{example}\label{ex:43}
Let $\la = (4,3)$ and
\[
  T \;=\; \begin{ytableau} 1 & 3 & 5 & 7 \\ 2 & 4 & 6 \end{ytableau}
  \quad\text{(in row form: $((1,3,5,7),\,(2,4,6))$)}.
\]
\textbf{Computation.} Numerically $p_T(11) = 0$.

\textbf{SC pairs.} Cell positions are $1 \mapsto (1,1)$, $2 \mapsto (2,1)$,
$3 \mapsto (1,2)$, $4 \mapsto (2,2)$, $5 \mapsto (1,3)$, $6 \mapsto (2,3)$,
$7 \mapsto (1,4)$. Adjacencies share a column iff $i$ is odd and $i < 7$:
SC pairs at $i \in \{1, 3, 5\}$.

\textbf{SR pairs.} No two consecutive integers lie in the same row.
SR pairs $= \emptyset$.

So $T$ has $p_T = 0$ and is rich in SC pairs but contains no SR pair.
The SR-companion conjecture fails on a $7$-cell tableau.
\end{example}

\subsection{An infinite family of counterexamples}

The same construction generalises: for any two-row partition
$\la = (k, m)$ with $k \ge m \ge 2$ and $k + m$ \emph{odd}, the tableau
\[
  T_{\la}^{\text{alt}} \;=\;
  \big(\,(1, 3, 5, \dots, k + m - 1),\; (2, 4, \dots, k + m - 2)\,\big)
\]
when $k = (k+m+1)/2$ \dots actually it is simplest to describe in
column-major form: $T$ is the unique tableau where the columns are
filled $(1,2), (3,4), (5,6), \dots$ left-to-right, with column-$k$
containing a single entry $k+m-1$ when $k > m$. Every adjacency is
a column descent; no adjacency is a same-row ascent. SC pair set
$= \{1, 3, 5, \dots\}$; SR pair set $= \emptyset$.

A direct calculation (omitted) shows $p_T = 0$ on at least the
sub-family $\la = (k, k-1)$ for $k \in \{3, 4, 5\}$. The asymptotic
behaviour is folded into the counts of Table~\ref{tab:results}.

\section{Where the symmetric argument fails}

\subsection{The asymmetry of SC and SR}

The combinatorics of $\SYT$ are not symmetric in row vs.\ column the way
this conjecture imagined. For a tableau $T$:
\begin{itemize}[leftmargin=2em]
\item SC pairs $= \mathrm{Des}(T)$, the descents.
\item SR pairs $= \mathrm{Asc}^{\!=}(T)$, the ``same-row ascents''
  (ascents that are not strict column moves).
\end{itemize}
Descents and same-row ascents are \emph{not} dual --- the dual of a
descent under the involution $T \mapsto T^t$ (transpose) is an
\emph{ascent}, but the ascent need not be a same-row one. A tableau
with rows shorter than its number of columns can have many descents
and zero same-row ascents (Example~\ref{ex:43} is exactly this).

In particular, $T^t$ has $\la^t$ as its shape, and our setup constrains
$\la = (\hla, 1^q)$. Transposition does not preserve the constraint
``every row $\hla_r \ge 2$'' (column-of-length-one rows transpose to
spread one-cell rows). So even the abstract duality is not available
within the dichotomy class.

\subsection{The chain-product positivity gap}

The Hoefsmit positivity that fuels Theorem~\ref{thm:sc} is a property of
the operator $T_i + 1$ in the $b' = 1$ convention. The natural mirror
operator is $1 + T_i$ acting on the right, which has the same matrix.
The rescaling that takes $b' = 1$ to $b = 1$ multiplies each $T_i + 1$
by a positive diagonal conjugation (axial-distance ratios), so each
\emph{individual} factor stays positive --- but the chain product
$L_n L_{n-1} \cdots L_{\ell+1}$ rearranges the order of the $S_i$.
Positivity of each factor is preserved by conjugation; positivity of a
specific \emph{matrix-product diagonal entry along the all-diagonal
path} is \emph{not}, because the conjugation rescales basis vectors
between factors, and the cancellations the all-diagonal rigidity rests
on do not survive.

Concretely, in Example~\ref{ex:43}, the matrix product
$\Omstar_{T,T}$ is also zero (this can be checked directly), but the
\emph{all-diagonal path} of $\Omstar$ at $T$ contains nonzero
contributions cancelled by an off-diagonal contribution of the opposite
sign --- precisely what positivity rules out for $\Om$. The
``positivity'' under $b = 1$ rescaling is in a different basis where
the path contributions for $\Omstar$ are not individually
nonneg-rational; rather, the rescaled matrix has nonneg-rational
entries, but the matrix-product expansion involves cross-terms from
different paths that are mutually cancellable.

\subsection{Strategic implication}

Diagnostic~3 (\texttt{2026-05-18-sr-companion.tex} for the SR forward
direction, \texttt{2026-05-17-leftmost-factor-per-syt-vanishing.tex}
for SC) stands in its current form: $p_T = 0$ implies SC pair \emph{or}
SR pair, with SC necessity universal and SR necessity local to the
projector-identity setting. They are not simultaneously necessary.

The converse of Diagnostic~4 cannot be reached by adapting
Theorem~\ref{thm:sc} symmetrically to $\Omstar$. Three structural routes
to the converse failed yesterday
(\texttt{2026-05-18-converse-diagnostic-4.tex}); today's negative result
rules out a fourth (``symmetric positivity rigidity''). The remaining
plausible routes are:
\begin{enumerate}[leftmargin=2em]
\item \textbf{Hecke-algebraic stratification.} Witness each zero $p_T$
  by a specific Hecke-algebraic identity in the chain product
  $R'_{\ell+1} \cdots R'_n$ (Conjectures~1 and~2 of the converse paper).
\item \textbf{Geometric / categorical lift.} The intrinsic Serre-pair
  structure on the braided $(\infty,2)$-cat of Soergel bimodules
  (LMRSW~2401.02956) supplies the right ``one of two pairings vanishes''
  fork at the level of $K$-theoretic classes on $\mathrm{Hilb}_n(\mathbb{C}^2)$.
\end{enumerate}
Today's refutation tightens the search: the structural input for the
converse is \emph{not} a positivity statement about $\Omstar$. It must
be specific to the algebraic structure of chain products.

\section{What this leaves open}

\begin{itemize}[leftmargin=2em]
\item \textbf{Refined Diagnostic~3.} Can the dichotomy ``SC or SR''
  at every $p_T = 0$ be strengthened to ``SC, with SR exactly when a
  specific structural condition $\mathcal{C}(\la, T)$ holds''? The
  counterexamples have visible pattern (interleaved two-row tableaux);
  partial converses on subfamilies are plausible.
\item \textbf{$\Omstar$-side per-SYT counterpart.} Define
  $p^*_T(q) := \Omstar_{T,T}$. Is $p^*_T = 0 \iff T$ has SR pair?
  This is the right symmetric statement, and a direct test on the
  $12$-shape corpus is one Python script away.
\item \textbf{Categorical bidegree refinement.} If the categorical
  lift produces an explicit bidegree obstruction for ``$\Om v = 0$ or
  $\Omstar v = 0$'', the SC/SR asymmetry may sharpen into a homological
  asymmetry between the two ``side''-bidegrees.
\end{itemize}

\section{Reproducibility}

The computational data underlying Table~\ref{tab:results} is generated
by \texttt{sr\_conjecture\_test.py} in
\texttt{\textasciitilde/projects/scratch/2026-05-18-converse-probe/},
runtime under one second. The script reuses the Hoefsmit seminormal
infrastructure built up over the preceding ten days
(\texttt{build\_T\_i\_seminormal} from
\texttt{\textasciitilde/projects/scratch/2026-05-08-rank-Pi/}). Each
zero diagonal is identified by exact comparison of $\Om_{T,T}$ to $0$
in $\QQ$ at $q = 11$; the integer-coefficient polynomials in $q$ that
vanish at $q = 11$ for the verified shapes also vanish symbolically (as
verified independently on smaller shapes in prior write-ups).

\end{document}
